\documentclass[12pt,a4paper]{article}
\usepackage[T1]{fontenc}
\usepackage[utf8]{inputenc}
\usepackage[spanish]{babel}
\usepackage{geometry}
\usepackage{booktabs}
\usepackage{xcolor}
\usepackage{listings}
\usepackage{hyperref}
\usepackage{enumitem}
\usepackage{titlesec}
\usepackage{parskip}

\geometry{margin=2.5cm}

\hypersetup{
  colorlinks=true,
  linkcolor=black,
  urlcolor=blue!70!black,
}

\lstdefinestyle{bash}{
  language=bash,
  basicstyle=\small\ttfamily,
  backgroundcolor=\color{gray!10},
  frame=single,
  rulecolor=\color{gray!40},
  breaklines=true,
  showstringspaces=false,
}

\title{%
  \textbf{Control de Iluminación ROG}\\[0.4em]
  \large Configuración personalizada en Linux para el\\
  ASUS ROG Zephyrus G14 GA403UV
}
\author{pc-config}
\date{\today}

\begin{document}

\maketitle
\tableofcontents
\newpage

% ─────────────────────────────────────────────
\section{Introducción}

Este documento describe la configuración de iluminación del teclado y el slash
bar (barra LED de la tapa) del portátil \textbf{ASUS ROG Zephyrus G14 GA403UV}
bajo un sistema \textbf{CachyOS} con el compositor \textbf{Hyprland} (entorno
Omarchy).

El objetivo es proporcionar tres funcionalidades:

\begin{enumerate}
  \item \textbf{Ciclar efectos Aura} del teclado con una tecla dedicada (Fn+F4).
  \item \textbf{Modo ambiental} (\emph{Screen Sync}): el teclado adopta
        dinámicamente el color dominante de la pantalla.
  \item \textbf{Ciclar modos del slash bar} con Super+\textbackslash.
\end{enumerate}

% ─────────────────────────────────────────────
\section{Pila de software}

\begin{center}
\begin{tabular}{lll}
  \toprule
  Componente & Versión & Función \\
  \midrule
  \texttt{asusctl} & 6.5.0-1 & CLI/daemon para hardware ASUS ROG \\
  \texttt{asusd}   & 6.5.0-1 & Servicio D-Bus (\texttt{xyz.ljones.Asusd}) \\
  \texttt{grim}    & ---      & Captura de pantalla en Wayland \\
  \texttt{python-pillow} & 12.2.0-1 & Reescalado de imagen para color dominante \\
  \texttt{busctl}  & systemd  & Lectura/escritura de propiedades D-Bus \\
  \texttt{notify-send} & libnotify & Notificaciones de escritorio \\
  \bottomrule
\end{tabular}
\end{center}

El daemon \texttt{asusd} expone los siguientes objetos D-Bus relevantes:

\begin{itemize}[nosep]
  \item \textbf{Teclado Aura}:\\
        \texttt{xyz.ljones.Asusd /xyz/ljones/aura/19b6\_3\_4 xyz.ljones.Aura}
  \item \textbf{Slash bar}:\\
        \texttt{xyz.ljones.Asusd /xyz/ljones/aura/193b\_4\_5 xyz.ljones.Slash}
\end{itemize}

% ─────────────────────────────────────────────
\section{Scripts instalados}

Todos los ejecutables se instalan en \texttt{\textasciitilde/.local/bin/} y
requieren únicamente permisos de usuario, a excepción de
\texttt{keyboard-aura-boot} (ver sección \ref{sec:arranque}), que se instala
como servicio de sistema con \texttt{setup-keyboard-aura-boot.sh}.

\subsection{\texttt{keyboard-ambient}}

Demonio Python que captura continuamente la pantalla con \texttt{grim}, calcula
el color dominante reescalando la imagen a 1×1 píxel, aplica un aumento de
brillo si el pico es inferior a 80, y envía el color al teclado vía:

\begin{lstlisting}[style=bash]
asusctl aura effect static -c RRGGBB
\end{lstlisting}

Parámetros opcionales:

\begin{center}
\begin{tabular}{lll}
  \toprule
  Argumento & Rango & Defecto \\
  \midrule
  \texttt{--fps} & 0.5\,--\,10 & 4 \\
  \texttt{--brightness} & 0.0\,--\,1.0 & 1.0 \\
  \bottomrule
\end{tabular}
\end{center}

\subsection{\texttt{keyboard-ambient-toggle}}

Script Bash que activa o desactiva el modo ambiental:

\begin{itemize}[nosep]
  \item \textbf{Al activar}: guarda el efecto Aura actual (propiedad
        \texttt{LedModeData} de D-Bus) en \texttt{/tmp/keyboard-aura-state} y
        lanza \texttt{keyboard-ambient} en segundo plano.
  \item \textbf{Al desactivar}: detiene el demonio y restaura el efecto previo
        reconstruyendo el comando \texttt{asusctl aura effect \ldots}\ a partir
        de los datos guardados.
\end{itemize}

\subsection{\texttt{keyboard-aura-cycle}}

Avanza al siguiente efecto Aura (\texttt{asusctl aura effect --next-mode}),
detiene el modo ambiental si estuviera activo, y muestra una notificación con
el nombre del efecto.

\textbf{Efectos soportados en el GA403UV} (número D-Bus → nombre):

\begin{center}
\begin{tabular}{rll}
  \toprule
  \# & Nombre & Descripción \\
  \midrule
  0  & Static           & Color fijo \\
  1  & Breathe          & Pulsación suave \\
  2  & Rainbow Cycle    & Ciclo de colores \\
  3  & Rainbow Wave     & Ola de arcoíris \\
  10 & Pulse            & Pulso \\
  \bottomrule
\end{tabular}
\end{center}

\texttt{asusd} expone únicamente estos cinco modos para este modelo; el resto de
efectos de la familia ROG no están disponibles (la detección los descarta).

\subsection{\texttt{keyboard-slash-cycle}}

Cicla entre los 16 modos del slash bar y muestra una notificación:

\begin{center}
\begin{tabular}{llll}
  \toprule
  Static & Bounce & Slash & Loading \\
  BitStream & Transmission & Flow & Flux \\
  Phantom & Spectrum & Hazard & Interfacing \\
  Ramp & GameOver & Start & Buzzer \\
  \bottomrule
\end{tabular}
\end{center}

\textbf{Nota importante:} El slash bar del modelo GA403UV solo se activa
durante eventos del sistema (arranque, suspensión y apagado). No existe modo
continuo; el cambio de modo se aplicará en el próximo evento.

\subsection{\texttt{keyboard-aura-boot}}

Aplica el efecto Aura y el brillo al arranque del sistema. No vive en
\texttt{\textasciitilde/.local/bin/} sino en \texttt{/usr/local/bin/}, y se
ejecuta mediante el servicio
\texttt{/etc/systemd/system/keyboard-aura-boot.service}. Ver la sección
\ref{sec:arranque} para el detalle y el motivo.

% ─────────────────────────────────────────────
\section{Efecto al arranque (Rainbow Cycle)}
\label{sec:arranque}

\subsection*{Por qué el teclado no arranca en Rainbow Cycle}

En Windows el efecto Aura lo escribe Armoury Crate/Aura Sync en el EC del
teclado y el EC lo conserva al apagar, por lo que el teclado se ilumina apenas
se da energía. En Linux eso no lo hace el firmware sino \texttt{asusd}, que
guarda el último estado en \texttt{/etc/asusd/aura\_<id>.ron} y lo reaplica
unos segundos después de arrancar el sistema (nunca durante BIOS/Limine).

Ese estado guardado no es necesariamente Rainbow Cycle:

\begin{itemize}[nosep]
  \item el modo cambia con \texttt{Fn+F4} (\texttt{keyboard-aura-cycle}) y con
        el modo ambiental (\texttt{Super+F4});
  \item \texttt{asusd} 6.5.0 reaplica al arrancar el modo y los power states,
        pero \textbf{no} el brillo: si quedó en Off, el teclado arranca
        apagado aunque el modo guardado sea Rainbow Cycle.
\end{itemize}

\subsection*{Solución instalada}

\texttt{setup-keyboard-aura-boot.sh} (paso 17) instala:

\begin{itemize}[nosep]
  \item \texttt{/usr/local/bin/keyboard-aura-boot} — aplica
        \texttt{asusctl aura effect rainbow-cycle} y
        \texttt{asusctl leds set med}, reintentando hasta 10 s si \texttt{asusd}
        todavía no responde.
  \item \texttt{keyboard-aura-boot.service} — \texttt{Type=oneshot},
        \texttt{After}/\texttt{Requires} de \texttt{asusd.service} y
        \texttt{Before=display-manager.service}, para que el efecto quede
        aplicado antes del login.
\end{itemize}

Para cambiar efecto o brillo sin tocar el script:

\begin{lstlisting}[style=bash]
sudo systemctl edit keyboard-aura-boot
# [Service]
# Environment=KEYBOARD_AURA_MODE=rainbow-wave
# Environment=KEYBOARD_AURA_BRIGHTNESS=high
\end{lstlisting}

\textbf{Límite de firmware:} durante POST/BIOS/Limine el teclado lo controla el
EC y el sistema operativo todavía no puede escribirle el efecto. El servicio
aplica justo antes del login; iluminarlo antes no es posible desde Linux.

% ─────────────────────────────────────────────
\section{Atajos de teclado}

Los atajos se registran en \texttt{\textasciitilde/.config/hypr/bindings.conf}
(Hyprland).

\begin{center}
\begin{tabular}{lll}
  \toprule
  Tecla & Acción & Script \\
  \midrule
  \textbf{Fn+F4} (XF86Launch3)     & Siguiente efecto Aura   & \texttt{keyboard-aura-cycle} \\
  \textbf{Super+F4}                 & Activar/desactivar modo ambiental & \texttt{keyboard-ambient-toggle} \\
  \textbf{Super+\textbackslash}     & Siguiente modo slash bar & \texttt{keyboard-slash-cycle} \\
  \bottomrule
\end{tabular}
\end{center}

\textbf{Nota sobre Fn+F4:} En este portátil la tecla Fn+F4 (AURA) genera el
keycode 202 (\texttt{KEY\_PROG3} = \texttt{XF86Launch3}), \emph{no}
\texttt{XF86KbdLightOnOff} como podría esperarse. Se detectó leyendo
directamente \texttt{/dev/input/event3}.

% ─────────────────────────────────────────────
\section{Instalación}

El script \texttt{setup-keyboard-ambient.sh} instala todo de forma idempotente:

\begin{lstlisting}[style=bash]
cd ~/Documentos/pc-config
bash setup-keyboard-ambient.sh
\end{lstlisting}

El script realiza las siguientes acciones:
\begin{enumerate}
  \item Escribe los cuatro scripts en \texttt{\textasciitilde/.local/bin/} con
        permisos de ejecución.
  \item Añade las tres líneas \texttt{bindd} a
        \texttt{\textasciitilde/.config/hypr/bindings.conf} si no existen.
\end{enumerate}

Para que el teclado arranque siempre en Rainbow Cycle, ejecutar además el paso
17 (requiere root):

\begin{lstlisting}[style=bash]
cd ~/Documentos/pc-config
sudo bash setup-keyboard-aura-boot.sh
\end{lstlisting}

Después de ejecutarlo, recarga Hyprland para que los atajos estén activos:

\begin{lstlisting}[style=bash]
hyprctl reload
\end{lstlisting}

% ─────────────────────────────────────────────
\section{Resolución de problemas}

\subsection*{El modo ambiental no restaura el efecto anterior}

Comprueba que el archivo de estado fue creado correctamente:

\begin{lstlisting}[style=bash]
cat /tmp/keyboard-aura-state
\end{lstlisting}

Si está vacío o no existe, verifica que \texttt{asusd} esté en ejecución:

\begin{lstlisting}[style=bash]
systemctl status asusd
\end{lstlisting}

\subsection*{Fn+F4 no responde}

Verifica el keycode real generado por esa tecla:

\begin{lstlisting}[style=bash]
python3 -c "
import struct, time, select
fmt = 'llHHI'; sz = struct.calcsize(fmt)
with open('/dev/input/event3', 'rb') as f:
    r,_,_ = select.select([f],[],[],5)
    if r:
        _,_,t,c,v = struct.unpack(fmt, f.read(sz))
        print(f'type={t} code={c} val={v}')
"
\end{lstlisting}

\subsection*{El teclado arranca apagado (o en otro efecto)}

Es el comportamiento por defecto de Linux: \texttt{asusd} reaplica el último
modo guardado y no restaura el brillo. Instalar el paso 17
(\texttt{setup-keyboard-aura-boot.sh}) y comprobar que el servicio esté activo:

\begin{lstlisting}[style=bash]
systemctl status keyboard-aura-boot.service
asusctl aura effect --next-mode   # revierte el modo ambiental
asusctl leds get                  # el brillo no debe estar en Off
\end{lstlisting}

Recordar que durante BIOS/Limine el teclado lo controla el EC y no se puede
iluminar desde Linux; el efecto se aplica justo antes del login.

\subsection*{El slash bar no cambia visualmente}

Es el comportamiento esperado en el modelo GA403UV. El slash bar solo muestra
animaciones durante eventos de sistema. El modo guardado se aplicará la próxima
vez que el equipo arranque, se suspenda o se apague.

\end{document}
